\chapter{Probability Refresher}\label{ch:probability}

Almost every definition and proof in cryptography is a statement about probabilities: ``the adversary guesses with probability $\frac{1}{2}$'', ``the ciphertext is independent of the message'', and so on. This appendix collects the few tools we need. Everything here is \textbf{discrete}: all our sets (keys, messages, ciphertexts) are finite.

\section{Sample space, events and probability}

Think of a random experiment, for example tossing a coin or choosing a key at random.
\begin{itemize}
  \item \textbf{Sample space} $\Omega$ $\to$ the set of all possible outcomes of the experiment (for example, all the keys in $\Kspace$);
  \item \textbf{Event} $E \subseteq \Omega$ $\to$ a set of outcomes we are interested in (for example, ``the first bit of the key is $0$'');
  \item \textbf{Probability} $\Pr$ $\to$ a function that assigns to every event a number between $0$ and $1$.
\end{itemize}

\begin{frameddefn}{Probability distribution (discrete)}
  A probability distribution on a finite set $\Omega$ assigns to every outcome $\omega$ a number $\Pr[\omega] \geq 0$ such that
  \[
    \sum_{\omega \in \Omega} \Pr[\omega] = 1
  \]
  The probability of an event is the sum of the probabilities of its outcomes: $\Pr[E] = \sum_{\omega \in E} \Pr[\omega]$.
\end{frameddefn}

Basic rules that follow immediately:
\begin{itemize}
  \item \textbf{Complement} $\to$ $\Pr[\overline{E}] = 1 - \Pr[E]$;
  \item \textbf{Disjoint events} $\to$ if $E_1 \cap E_2 = \emptyset$, then $\Pr[E_1 \cup E_2] = \Pr[E_1] + \Pr[E_2]$;
  \item \textbf{General union} $\to$ $\Pr[E_1 \cup E_2] = \Pr[E_1] + \Pr[E_2] - \Pr[E_1 \cap E_2]$.
\end{itemize}

\section{The uniform distribution}

The most important distribution in cryptography: \textbf{keys are always chosen uniformly}.

\begin{frameddefn}{Uniform distribution}
  The uniform distribution on a finite set $S$ gives the same probability to every element:
  \[
    \Pr[x] = \frac{1}{|S|} \qquad \text{for every } x \in S
  \]
  Notation: $x \sample S$ means ``$x$ is chosen uniformly at random from $S$''.
\end{frameddefn}

For example, with $K \sample \{0,1\}^n$ there are $2^n$ keys, so every key has probability $2^{-n}$. The probability of an event then becomes simple \textbf{counting}, $\Pr[E] = \frac{|E|}{|S|}$: the probability that the first bit of the key is $0$ is $\frac{2^{n-1}}{2^n} = \frac{1}{2}$.

\section{Random variables}

\begin{frameddefn}{Random variable}
  A random variable $X$ is a function that assigns a value to every outcome of the experiment, $X : \Omega \to V$. The \textbf{distribution} of $X$ is the list of probabilities $\Pr[X = v]$ for every value $v \in V$.
\end{frameddefn}

In our setting there are three random variables:
\begin{itemize}
  \item \textbf{the key} $K$ $\to$ uniform over $\Kspace$;
  \item \textbf{the message} $M$ $\to$ any distribution over $\Mspace$ (what Eve knows a priori), independent of $K$;
  \item \textbf{the ciphertext} $C = \Enc(K, M)$ $\to$ its distribution is determined by those of $K$ and $M$.
\end{itemize}

\paragraph{The subscript notation.} $\Pr_K[\Enc(K, m) = c]$ means that the message $m$ is \textbf{fixed}, and the probability is computed \textbf{only over the random choice of the key} $K$. The subscript tells you which randomness is involved.

\section{Joint and conditional probability}

Often we want to talk about two random variables at the same time, or about one of them \textbf{after} learning something about the other.
\begin{itemize}
  \item \textbf{Joint probability} $\to$ $\Pr[X = x \wedge Y = y]$, also written $\Pr[X = x, Y = y]$, is the probability that \textbf{both} things happen;
  \item \textbf{Conditional probability} $\to$ $\Pr[X = x \mid Y = y]$ is the probability that $X = x$ \textbf{knowing that} $Y = y$ happened.
\end{itemize}

\begin{frameddefn}{Conditional probability}
  If $\Pr[Y = y] > 0$:
  \[
    \Pr[X = x \mid Y = y] = \frac{\Pr[X = x \wedge Y = y]}{\Pr[Y = y]}
  \]
\end{frameddefn}

\textbf{Intuition.} Knowing that $Y = y$ happened \textbf{restricts the world} to the outcomes where $Y = y$. Inside this smaller world, we measure how often $X = x$ also happens.

The two rules we use most often, both obtained by rearranging the definition:
\begin{itemize}
  \item \textbf{Multiplication rule} $\to$ $\Pr[X = x \wedge Y = y] = \Pr[X = x \mid Y = y] \cdot \Pr[Y = y]$;
  \item \textbf{Symmetry} $\to$ the joint probability can be split in both directions: $\Pr[X = x \mid Y = y] \cdot \Pr[Y = y] = \Pr[Y = y \mid X = x] \cdot \Pr[X = x]$.
\end{itemize}

\section{Independence}

\begin{frameddefn}{Independent random variables}
  $X$ and $Y$ are \textbf{independent} if for every $x$ and $y$:
  \[
    \Pr[X = x \wedge Y = y] = \Pr[X = x] \cdot \Pr[Y = y]
  \]
  Equivalently (when $\Pr[Y = y] > 0$), $\Pr[X = x \mid Y = y] = \Pr[X = x]$.
\end{frameddefn}

\textbf{Intuition.} Learning the value of one variable does not change the distribution of the other. Two different coins tossed separately are independent; the first bit of a key and the whole key are not. In our setting the key $K$ is always independent of the message $M$, and statement (ii) of \Cref{lem:equivalent} says exactly that $M$ and $C$ are independent.

\section{Law of total probability and Bayes' rule}

\begin{framedthm}{Law of total probability}
  For any random variables $X$ and $Y$:
  \[
    \Pr[X = x] = \sum_{y} \Pr[X = x \mid Y = y] \cdot \Pr[Y = y]
  \]
  where the sum is over the values $y$ with $\Pr[Y = y] > 0$.
\end{framedthm}

\textbf{Intuition.} To compute the probability of $X = x$, we split the world into cases according to the value of $Y$, compute the probability in each case, and take the \textbf{weighted average}, where each case weighs as much as its probability. If every case gives the same value, the average is that value too: this is exactly step (f) in the proof of (iii) $\Rightarrow$ (i).

\begin{framedthm}{Bayes' rule}
  If $\Pr[Y = y] > 0$:
  \[
    \Pr[X = x \mid Y = y] = \frac{\Pr[Y = y \mid X = x] \cdot \Pr[X = x]}{\Pr[Y = y]}
  \]
\end{framedthm}

\textbf{Intuition.} Bayes' rule \textbf{reverses the conditioning}. In cryptography it is natural to compute $\Pr[C = c \mid M = m]$ (``if the message is $m$, how likely is the ciphertext $c$?''), which depends only on the key. But what Eve cares about is the opposite, $\Pr[M = m \mid C = c]$ (``I see $c$: how likely is the message $m$?''), and Bayes connects the two. $\Pr[X = x]$ is called the \textit{prior} (a priori) and $\Pr[X = x \mid Y = y]$ the \textit{posterior} (a posteriori).

\section{The union bound}

\begin{framedthm}{Union bound}
  For any events $E_1, \ldots, E_t$ (not necessarily disjoint):
  \[
    \Pr[E_1 \cup E_2 \cup \cdots \cup E_t] \leq \sum_{i=1}^{t} \Pr[E_i]
  \]
\end{framedthm}

\textbf{Intuition.} The probability that \textbf{at least one} of several bad things happens is at most the sum of their probabilities: when the events overlap, we are just counting some outcomes more than once. It appears everywhere in computational security: if each of polynomially many bad events has negligible probability, then the probability that at least one of them happens is still negligible.

\section{Expectation}

\begin{frameddefn}{Expected value}
  The expected value (average) of a random variable $X$ with numerical values is
  \[
    \E[X] = \sum_{x} x \cdot \Pr[X = x]
  \]
\end{frameddefn}

\textbf{Linearity of expectation}: $\E[X + Y] = \E[X] + \E[Y]$ and $\E[aX] = a\,\E[X]$, \textbf{even if $X$ and $Y$ are not independent}. For example, if $X$ is the number of heads in $n$ tosses of a fair coin, writing $X$ as the sum of $n$ variables that are $1$ on heads and $0$ on tails (each with expectation $\frac{1}{2}$) gives $\E[X] = \frac{n}{2}$ with no further computation.

\section{A key fact: XOR with a uniform key}

This small lemma is the heart of the One-Time Pad, and it comes back many times in the course.

\begin{framedlem}{XOR with a uniform string}
  Let $K \sample \{0,1\}^n$ and let $M$ be \textbf{any} random variable over $\{0,1\}^n$, \textbf{independent} of $K$. Then $K \oplus M$ is uniform over $\{0,1\}^n$.
\end{framedlem}

\begin{proof}
  Fix any $c \in \{0,1\}^n$. By the law of total probability (splitting by the value of $M$) and independence:
  \[
    \Pr[K \oplus M = c] = \sum_{m} \Pr[M = m] \cdot \Pr[K = m \oplus c] = \sum_{m} \Pr[M = m] \cdot 2^{-n} = 2^{-n}
  \]
  The result is $2^{-n}$ for every $c$, so $K \oplus M$ is uniform.
\end{proof}

\textbf{Intuition.} Whatever is XOR-ed with a uniform, independent key becomes uniform noise. For a single bit: if the key bit is a fair coin, flipping or not flipping the message bit according to it gives a fair coin, whatever the message bit was.
